%% ICT Express submission — revised manuscript
%% Format: Elsevier CAS double-column (cas-dc), per the current official
%% Guide for Authors (els-cas-templates is the template bundle the journal
%% currently recommends; LaTeX two-column is permitted for LaTeX submissions).
%% All numbers in tables/figures are \input from script-generated files in
%% ../tables and ../figures; see ../traceability.csv for provenance.

\documentclass[a4paper,fleqn]{cas-dc}

\usepackage[numbers,sort&compress]{natbib}
\usepackage{booktabs}
\usepackage{graphicx}
\usepackage{amsmath}
\usepackage{url}
\usepackage{caption}

\newcommand{\ac}[0]{\ensuremath{ac_1}}

\begin{document}
\let\WriteBookmarks\relax
\def\floatpagepagefraction{1}
\def\textpagefraction{.001}

\shorttitle{What does the LLM actually contribute?}
\shortauthors{Y. Sun and N. Yin}

\title[mode = title]{What Does the LLM Actually Contribute? An Evidence-Constrained
Audit of LLM-Guided Retrieval over Heterogeneous Bilingual RDF Graphs}

\author[1]{Yuxiang Sun}[type=editor,
  orcid=0000-0003-0165-7664]
\credit{Conceptualization, Methodology, Software, Investigation, Writing -- original draft}
\ead{syx921120@gmail.com}
\affiliation[1]{organization={Guangzhou University of Software},
            city={Guangzhou},
            country={China}}

\author[2]{Nan Yin}[orcid=0009-0007-3506-9592]
\credit{Validation, Data curation, Writing -- review \& editing}
\ead{ny@caist.ca}
\affiliation[2]{organization={Canada Advanced Institute of Science and Technology},
            city={Ottawa},
            country={Canada}}

\cortext[1]{Corresponding author.}

\begin{abstract}
Retrieval pipelines that place a large language model (LLM) inside
knowledge-graph retrieval are usually evaluated end-to-end, making gains hard
to attribute. We re-examine one such pipeline -- LLM-guided alignment plus
bounded traversal over two heterogeneous bilingual RDF graphs (no
\texttt{owl:sameAs}) -- under three controls. First, an evidence-constrained
metric (\ac{}) replaces a judge-proxy metric coupled to the answer generator.
Second, an alignment-supervision probe re-derives the planner's prompt: on a
49{,}781-triple, 500-query synthetic benchmark the gold cross-lingual URI is
present in 130/150 queries; Hits@1 is 0.500 when present versus 0/20 observed
successes when absent (exact 95\% CI $[0,0.168]$). Third, real SPARQL~1.1 and
multilingual dense baselines replace capability-limited surrogates, and a
matched-supervision control gives the compared systems the same alignment
table. Under that control a deterministic linker matches the 7B planner exactly
on cross-lingual \ac{} (0.433; zero discordant pairs), so the original
advantage reflected unequal supervision. The LLM retains an overall planning
advantage (0.772 vs.\ 0.498; $p<10^{-20}$), but cross-lingual ability comes
from the supplied table, not autonomous alignment. Real SPARQL restores
two-hop accuracy to 0.900; localized traversal cuts logical triple access by
$\sim$3 orders of magnitude at unchanged accuracy, while the LLM dominates
latency. A small authentic Wikidata check supports methodology transfer but
not cross-graph alignment. Code, data, and per-query records are released.
\end{abstract}

\begin{keywords}
knowledge graph \sep RDF retrieval \sep cross-lingual entity alignment \sep large language model \sep evaluation methodology
\end{keywords}

\maketitle

\section{Introduction}\label{sec:intro}
GraphRAG-style pipelines increasingly put an LLM inside the retrieval loop: the
model reads the question, emits a plan (aligned entity URIs, predicate
constraints), and a traversal engine executes the plan over a store. Reported
end-to-end gains are easy to produce and hard to interpret, for three reasons.
The answer generator and the judge may be the same model (inflating every
metric); baselines may be weaker than the real methods they claim to represent;
and when the LLM's prompt already contains ground-truth alignment information,
``the LLM aligned the entity'' may actually mean ``the planner copied a URI it
was given''~\cite{pan2024unifying,barnett2024seven}.
Prior work spans LLM--KG unification~\cite{pan2024unifying}, GraphRAG-style
graph access~\cite{edge2024graphrag}, cross-lingual entity
alignment~\cite{sun2018jape,yang2024segalign}, multilingual dense
retrieval~\cite{e5,zhang2023miracl}, and evidence-grounded
evaluation~\cite{thorne2018fever}. These strands typically report end-to-end
gains or optimise alignment itself; they do not ask whether LLM planning keeps
an advantage when the same alignment table, workload, and literal-evidence
criterion are given to a deterministic retriever. Our audit targets that
attribution gap rather than proposing a new architecture.

This letter is an evidence-constrained re-evaluation of such a pipeline for
retrieval over two heterogeneous bilingual RDF graphs. The setting is
deliberately synthetic: graph~A uses Chinese labels and schema~X predicates,
graph~B uses English labels and schema~Y predicates, the two graphs describe
the same persons under different URIs, and there is no cross-graph identity
link. We make three contributions:
\begin{enumerate}
\item \textbf{A judge-free, evidence-constrained metric.} \ac{} counts a query
as correct only if the gold entity is ranked first \emph{and} the gold value is
literally present in the retrieved evidence triples; $v_{ic}$ drops the ranking
requirement. Both are computed from saved per-query records and replace a
success metric that was coupled to a deterministic mock answerer.
\item \textbf{An alignment-supervision probe.} For every query we re-derive the
LLM planner's prompt and mark whether the gold URI appears in it verbatim. This
single measurement separates ``the model solved it'' from ``the model selected
what it was shown'' and, to our knowledge, is not routinely reported in
LLM-in-the-loop retrieval evaluations.
\item \textbf{Capability-corrected baselines, matched supervision, and cost
accounting.} We replace a ``SPARQL'' baseline that never executed SPARQL with
real SPARQL~1.1 via rdflib, and a monolingual hashed-encoder vector baseline
with multilingual dense retrieval (e5-small) plus traversal. Crucially, we also
re-run a deterministic localized linker with \emph{the same alignment table}
given to the planner, which isolates alignment information from LLM planning.
Finally we separate serial service time, loaded response time, and logical
triple access (GSR), an access ratio rather than I/O or memory.
\end{enumerate}
The headline outcome is a demotion, not a promotion: after the corrections,
the pipeline still leads on full-benchmark accuracy (\ac{} $0.772$ vs.\
$0.498$ for a non-LLM linker under the same alignment table and $0.646$ for
unaided real SPARQL; Holm-corrected $p<10^{-20}$ and $p<3.2\times10^{-8}$), but
that linker matches its cross-lingual \ac{} exactly ($0.433$; zero discordant
pairs in 150 queries). The surviving contribution is planning and predicate
selection on in-language queries, not cross-lingual alignment
(Sections~\ref{sec:probe}--\ref{sec:baselines}). A transliteration baseline
solves every two-character-name query perfectly at $\sim$0.1\,ms
(Section~\ref{sec:baselines}), and the LLM dominates end-to-end latency by
roughly five orders of magnitude over traversal. Although our benchmark is
synthetic, the audit protocol itself is domain-agnostic;
Section~\ref{sec:conclusion} draws out what it implies for three ICT settings
with the same structure: infrastructure metadata management, cross-domain
telecom knowledge-graph integration, and heterogeneous IoT retrieval.

\section{System and benchmark}\label{sec:system}
\textbf{Store.} A deterministic generator builds two rdflib \texttt{Graph}s
with a target of 50{,}000 triples (realized: 49{,}781 triples, 3{,}481
entities, 8{,}218 labels/aliases). Three noise sources are built in: polysemy
(same label, different identities), schema mismatch (synonymous predicates
within each schema), and a cross-lingual naming gap (labels share no
characters across graphs; for two-character names the pair is a romanisation
pair, for three-character names it is not). The generator, queries, seeds,
and all derived records are released with the letter (see the code and data
availability statement), so every number below can be recomputed from the
per-query CSVs alone.

\textbf{Pipeline.} A surface stage matches the query against the label/alias
index in both graphs; an LLM planner (qwen2.5:7b served by Ollama, temperature
0, JSON output, 256-token budget) receives the matched anchors together with
\emph{entity profiles} -- the canonical entity's surface forms in both
languages and all URIs denoting it -- and emits aligned URIs plus predicate
constraints; a localized executor then traverses only the aligned subjects'
triples, expanding one hop along work links when the intent targets work
attributes. The planner is wrapped with a simulated parametric memory whose
coverage (92\%) is an input parameter modelling the long tail an LLM has never
seen; coverage is disclosed rather than measured.

\textbf{Queries.} 500 queries in three families: 200 one-hop attribute
questions, 150 two-hop work-attribute questions, and 150 cross-lingual
questions whose gold counterpart lives only in the other graph. Ground truth
is recorded by the generator, so no answer is invented by a judge.

\textbf{Notation and metrics.} For query $q$, let $R(q)$ be the system's ranked
answer list with rank operator $\rho(q)\in\{1,2,\dots\}$ of the gold entity,
$E(q)\subseteq T$ the evidence triples returned to the answer stage, $v^{*}(q)$
the generator-recorded gold value, and $S(q)\subseteq T$ the triples the
executor logically scanned. Given $n$ queries, the ranking metrics are
\begin{equation}
\mathrm{Hits@1} \;=\; \frac{1}{n}\sum_{q}\mathbf{1}\!\left[\rho(q)=1\right],
\qquad
\mathrm{MRR} \;=\; \frac{1}{n}\sum_{q}\frac{1}{\rho(q)} .
\label{eq:rank}
\end{equation}
The evidence-constrained metrics couple ranking to grounding:
\begin{align}
\ac{}(q) &= \mathbf{1}\!\left[\rho(q)=1 \wedge v^{*}(q)\in E(q)\right],
\label{eq:ac1}\\
v_{ic}(q) &= \mathbf{1}\!\left[v^{*}(q)\in E(q)\right],
\label{eq:vic}
\end{align}
so \ac{} scores a query only when the gold entity is ranked first \emph{and}
the gold value is literally present in the returned evidence; $v_{ic}$ drops
the ranking conjunct. Neither can be satisfied by a judge agreeing with the
answer generator. Logical access is the ratio
\begin{equation}
\mathrm{GSR}(q)\;=\;\frac{|S(q)|}{|T|}\in[0,1],
\label{eq:gsr}
\end{equation}
reported strictly as an access ratio over the store $T$, not as measured
I/O or memory. Accuracy and GSR are aggregated over all 500 queries per
system.

\textbf{Paired tests.} For a binary metric, systems A and B are compared on
the identical query set with McNemar's exact test on the discordant pairs:
writing $b$ for the count with $ac_{1,A}(q)=1,\,ac_{1,B}(q)=0$ and $c$ for the
reverse,
\begin{equation}
p \;=\; 2\!\!\sum_{i=0}^{\min(b,c)}\!\!\binom{b+c}{i}\Big/2^{\,b+c},
\label{eq:mcnemar}
\end{equation}
the two-sided exact tail of $\mathrm{Bin}(b{+}c,\tfrac12)$. For the
continuous GSR we use the Wilcoxon signed-rank statistic on paired
differences $d_q=\mathrm{GSR}_A(q)-\mathrm{GSR}_B(q)$, and Holm's
step-down correction over the family of contrasts within each metric:
$p_{(k)}^{\mathrm{Holm}}=\max_{j\le k}\,(m{-}j{+}1)\,p_{(j)}$ for the ordered
$m$ raw $p$-values. For a binomial proportion $k/n$ we use the two-sided
Clopper--Pearson exact interval
\begin{equation}
\mathrm{CI}_{1-\alpha}=\left[B^{-1}\!\left(\tfrac{\alpha}{2};k,n-k+1\right),\;
B^{-1}\!\left(1-\tfrac{\alpha}{2};k+1,n-k\right)\right],
\label{eq:cp}
\end{equation}
where $B^{-1}(\cdot;a,b)$ is the Beta quantile; this avoids a normal
approximation when $k=0$.

\section{What the probe shows}\label{sec:probe}
The planner's prompt is built from a ground-truth surface-form table. When the
simulated memory covers an entity, the correct cross-lingual counterpart URI is
inside the prompt. Let $P(q)\subset T$ be the set of URIs that appear in the
planner's prompt for query $q$ (re-derived deterministically from the same
seeds as the main run), $u^{*}(q)$ the gold entity URI, and let
\begin{equation}
g(q)\;=\;\mathbf{1}\!\left[u^{*}(q)\in P(q)\right]
\label{eq:probe}
\end{equation}
mark the \emph{alignment-supervision} condition. The probe estimator simply
stratifies any metric $M$ by this condition,
$\hat{M}_{g=i}=\tfrac{1}{n_i}\sum_{q:\,g(q)=i}M(q)$ with $n_i$ the stratum
size, so a system can be credited with cross-lingual accuracy only on queries
where the gold URI was \emph{not} already shown to it.

Table~\ref{tab:probe} and Figure~\ref{fig:probe} give the outcome split. For
130 of the 150 cross-lingual
queries (86.7\%) the gold URI was in the prompt; the planner achieved Hits@1
$0.500$ on those (exact 95\% CI $[0.411,0.589]$) and 0/20 observed successes on
the uncovered queries (exact 95\% CI $[0,0.168]$ by Eq.~\eqref{eq:cp}). Thus no
success was observed without exposure, but zero observed successes do not
establish zero capability. The cross-lingual ability of the pipeline is
therefore bounded by the supplied alignment table, modulated by a
near-coin-flip selection rate; it is not evidence of model-supplied world
knowledge. One structural caveat: the two-hop family's gold answer is a
\emph{work} URI reached by traversal, so the probe does not apply there (marked
n/a); those queries succeed through person alignment plus the two-hop executor,
which the ablations confirm (removing the predicate constraint collapses two-hop
\ac{} from 0.933 to 0.000).

\input{../tables/table2_probe}
\captionof{table}{Alignment-supervision probe: outcomes stratified by whether
the gold URI was present in the planner's prompt.}\label{tab:probe}

\begin{figure}
\centering
\includegraphics[width=\columnwidth]{../figures/fig3_probe.pdf}
\caption{Alignment-supervision probe. When the gold URI is in the planner's
prompt the 7B planner selects it only half the time on cross-lingual queries;
when it is absent no success was observed in 20 queries, with an exact 95\%
interval reaching 0.168. On one-hop queries the prompt
almost always contains the gold entity (surface matching), and accuracy is
near-perfect.}
\label{fig:probe}
\end{figure}

\section{Capability-corrected baselines}\label{sec:baselines}
The original evaluation compared against (i) a ``SPARQL'' baseline that walked
Python triple patterns without a SPARQL engine and scored 0.000 on two-hop and
cross-lingual families, and (ii) a dense baseline with a monolingual hashed
encoder that scored 0.000 cross-lingually by construction. Neither number
supports a claim about SPARQL or dense retrieval in general.

Table~\ref{tab:main} and Figure~\ref{fig:acc} add three fair systems on the
identical store and query set. \emph{Real SPARQL~1.1} (rdflib engine; two-stage
resolution over \texttt{rdfs:label} and alias predicates, \texttt{VALUES}
push-down, two-hop work pattern) recovers two-hop \ac{} to $0.900$ -- the legacy
number of $0.000$ was an implementation artefact -- and reaches $0.940$ on
one-hop. It scores $0.000$ cross-lingually, which \emph{is} a genuine result:
without \texttt{owl:sameAs} no query can cross the naming gap unaided.
\emph{Multilingual dense retrieval} (e5-small, 8{,}218 passage documents, mean
pooled, L2-normalised) with traversal reaches $0.925$ on one-hop but $0.580$
on two-hop and $0.000$ cross-lingually; removing traversal collapses two-hop
to $0.000$, isolating what the graph step contributes. \emph{Dense without
traversal} reaches only $0.370$ overall: neither head has a cross-lingual
route, so this family of methods is structurally capped on the cross-lingual
third of the benchmark.

\emph{Transliteration.} A small pinyin matcher with identical traversal scores
$0.640/0.613$ on one-hop/two-hop and $0.307$ cross-lingually at $\sim$0.08\,ms
per query (mean, $n=198$). Diagnosis matters more than the headline: its
cross-lingual successes are exactly the 46 two-character-name queries (Hits@1
$1.000$, since those pairs are romanisation pairs by construction) and none of
the 104 three-character-name queries (Hits@1 $0.000$). The $0.307$ is a mixture
fixed by the generator's name-length distribution and must not be read as
evidence that phonetic matching is a strong cross-lingual linker.

\emph{Matched alignment supervision.} The original cross-lingual comparison
confounded capability with information: the planner's prompt contained a
truth-derived alias table, while SPARQL and dense retrieval did not. We
therefore rebuilt the planner's lookup path
(\texttt{knowledge}$\to$\texttt{label\_to\_key}$\to$\texttt{alias\_index}) and
ran a deterministic localized linker under three conditions: no alignment
table, the same matched table as the planner, and oracle gold URIs (an upper
bound only, not a deployable baseline).
Table~\ref{tab:matched} reports the cross-lingual subset. With the same table
the deterministic linker and the 7B planner are indistinguishable query by
query ($0.433$ each; 0/150 discordant pairs). Without the table both score
$0.000$; the oracle linker reaches $1.000$. The original cross-lingual lead is
therefore attributable to supplied alignment information rather than LLM
reasoning. On the full 500-query set the LLM still leads the matched linker
($0.772$ vs.\ $0.498$); its simplified two-hop ranker understates a fully tuned
non-LLM system, but this does not affect the cross-lingual rank-by-rank
equality.

For accuracy contrasts, Holm-adjusted $p<10^{-3}$ throughout ($n=500$ paired
queries); fair-baseline contrasts satisfy $p\le3.2\times10^{-8}$, whereas the
strongest ablation, \texttt{no\_aligner}, is $p=7.85\times10^{-4}$. The margin
over the strongest fair baseline is $0.772-0.646=0.126$, roughly half of what
the legacy baselines suggested. The in-language ablation also cuts against the
design: dropping the aligner entirely raises one-hop, two-hop and overall
\ac{} to $1.000$, because surface matching plus traversal already suffices
when no cross-graph alignment is required -- the aligner pays a tax
(mis-selection) on the two-thirds of queries that never needed it, connecting
to the known failure mode of retrieving the wrong entity before reasoning
begins \cite{barnett2024seven}.

\input{../tables/table1_main}
\captionof{table}{Main results at 49{,}781 triples / 500 queries. \ac{} is the
evidence-constrained strict metric; GSR is a logical access ratio, not measured
I/O. $^{\dagger}$Legacy baselines, shown for comparability only; claims rest on
the fair systems. MRR of the full system is 0.874.}\label{tab:main}

\input{../tables/table_matched}
\captionof{table}{Matched-supervision control on the 150 cross-lingual queries.
The matched condition gives the compared systems the same truth-derived alias
table; oracle rows are upper bounds only. The LLM and deterministic linker have
0/150 discordant pairs under matched supervision.}\label{tab:matched}

\begin{figure}
\centering
\includegraphics[width=\columnwidth]{../figures/fig2_accuracy.pdf}
\caption{Strict accuracy (\ac{}) per query family. The real SPARQL baseline
recovers two-hop accuracy that the legacy implementation reported as zero;
cross-lingual accuracy is available only to systems that are handed an
alignment table.}
\label{fig:acc}
\end{figure}

\section{Cost structure}\label{sec:cost}
Table~\ref{tab:latency} reports a unified 24-query workload: all executable
systems received the same qids, order, warm-up and cache policy. On this serial
service-time protocol the LLM planner costs $8.38$\,s per query on average
(median $8.58$\,s), of which $\sim$99.99\% is model time, while traversal and
non-LLM linkers are sub-millisecond. A separate four-worker run confirms that
queueing, not parallelism, dominates when one local model backend is shared;
loaded response time is therefore reported separately from service time. GSR
behaves as
designed: removing localized pruning leaves accuracy identical (\ac{} $0.772$,
$p=1$) while logical access jumps to $1.0$ of the store and traversal time
from $0.19$\,ms to $47.1$\,ms ($\sim$245$\times$), i.e.\ pruning is a pure
cost lever worth $\sim$3 orders of magnitude of logical access, with no
accuracy claim attached.
The scaling sweep confirms the same structure as the store grows. We fit the
log--log model
\begin{equation}
\log \overline{\mathrm{GSR}} \;=\; a + \beta\,\log|T|
\label{eq:fit}
\end{equation}
at sizes $|T|\in\{10^4, 5{\times}10^4, 10^5, 2{\times}10^5\}$: across
10K--200K triples the localized variants' GSR decays with $\beta=-0.98$ to
$-1.02$ ($R^2 \ge 0.999$), but this is a denominator effect: the mean number
of logically scanned triples stays 17.5--19.3 across the same 20-fold size
range (exponent $0.015$, $R^2=0.195$). We therefore do not claim overall
$O(1)$ complexity; the supported statement is that observed logical access
stays approximately constant under the evaluated scaling conditions. The
unpruned variant stays at $1.0$
and the SPARQL systems at $\approx 0.156$ of the store regardless of size
(Fig.~\ref{fig:scaling}). Real SPARQL latency grows with store size, whereas
the full system stays LLM-dominated. On the unified 24-query workload, real
SPARQL costs $290$\,ms mean / $157$\,ms median versus $2.6$\,ms for
vector-GraphRAG and $0.09$\,ms for the transliteration linker, i.e.\ one to
five orders of magnitude below the LLM. Dense e5 is excluded from this
unified latency run because torch/transformers were unavailable in that
environment; its accuracy remains reported on the 500-query benchmark.

\begin{figure}
\centering
\includegraphics[width=\columnwidth]{../figures/fig4_scaling.pdf}
\caption{GSR (logical access ratio) vs.\ store size, log--log. Localized
variants decay as $|T|^{-1}$; unpruned and SPARQL-family access is
size-proportional.}
\label{fig:scaling}
\end{figure}

\input{../tables/table3_latency}
\captionof{table}{Unified serial-service workload (same 24 qids, fixed order,
two warm-up queries, two timed repeats, plan cache cleared). Dense e5 was not
executable in this latency environment; its 500-query accuracy is reported
elsewhere. The $n$ column is the number of timed queries.}\label{tab:latency}

\section{Threats to validity and scope}\label{sec:limits}
\textbf{Synthetic main benchmark; limited real-RDF check.} System comparisons
come from one generated benchmark; there is no real ICT/telecommunications
deployment claim. A separate transfer check used 632 authentic Wikidata zh/en
triples (70 entities) and 184 offline queries (114 one-hop, 70 cross-lingual):
one-hop \ac{}=1.000 in all three supervision groups. Cross-lingual matched
\ac{} was 0.700 overall ($49/70$), exactly the train-pool alignment coverage;
the 21 held-out non-oracle queries scored 0/21 (exact 95\% CI $[0,0.161]$) and
the oracle scored 1.000. This supports metric and stratification transfer, not
cross-graph alignment: Wikidata is identity-linked, and two-hop, ambiguous,
and predicate-mapping families were not executed. The synthetic cross-lingual family tests whether an alignment
table was \emph{provided}, not whether a system can \emph{infer} alignment;
systems without such a table are structurally capped at zero here.
\textbf{Simulated memory.} The 92\% coverage parameter controls the
cross-lingual upper bound directly; it is a disclosed input, not a measured
model property. \textbf{One model, one machine.} Latency numbers are specific
to qwen2.5:7b on a local Ollama server and would change with a hosted API.
\textbf{No independent human annotation.} \ac{}/$v_{ic}$ are string-match
checks over saved evidence, not human judgements; the earlier mock-judge
metrics were removed precisely because they could be. \textbf{GSR} counts
logical triple access; it is not measured I/O, memory, or dollars.

\section{Conclusion}\label{sec:conclusion}
With a judge-free metric, a supervision probe, capability-corrected baselines,
and a matched-supervision control, the same benchmark tells a more defensible
story. Localized traversal is a pure cost lever: it cuts logical access by
$\sim$3 orders of magnitude without changing accuracy. The LLM's real gain is
planning and predicate selection on in-language queries ($0.772$ vs.\ $0.498$
for the matched linker), whereas cross-lingual \ac{} is identical to a
deterministic linker once both receive the same table ($0.433$). Answering
RQ1--RQ5: (i) the LLM beats unaided SPARQL/dense overall, but not the matched
linker on cross-lingual queries; (ii) cross-lingual performance is inherited
from the supplied table; (iii) the protocol transfers to authentic Wikidata
without establishing cross-graph alignment; (iv) the remaining contributions
decompose into alignment table, planner, and cost-only traversal; and (v) LLM
latency is justified only for planning-heavy queries, not for cross-lingual
linking. The actionable lessons are methodological: report what the prompt
contains, judge with metrics the answer generator cannot influence, and let
baselines use the capability they claim to represent.

\textbf{Implications for ICT systems.} The three controls transfer directly
to ICT settings that share the structure studied here -- heterogeneous
schemas, bilingual or multilingual labelling, and no shared identity links.
In \emph{ICT infrastructure metadata management}, asset and catalogue
metadata spread across CRs, inventory and topology systems is exactly the
``two schemas, one entity'' situation; before an LLM planner is credited
with ``resolving'' an asset, the same probe can ask whether the mapped
identifier was already in its prompt. In \emph{cross-domain telecom
knowledge-graph integration}, operator and vendor graphs maintained in
different languages will realistically be bridged by curated alignment
tables rather than \texttt{owl:sameAs}; Eq.~\eqref{eq:probe} gives
integrators a cheap audit of how much of the measured accuracy is inherited
from the table rather than produced by the model. In \emph{heterogeneous
IoT retrieval}, where device, platform, and application ontologies coexist,
the cost findings (LLM seconds vs.\ sub-millisecond traversal; near-constant
scanned-triple counts in the scaling sweep) bound where an in-the-loop planner
is affordable --
planning is justified for in-language predicate and plan selection, whereas
cross-graph alignment should be audited as an explicit information source
rather than credited to the model. Future work must move the evaluation
to settings where the alignment table is not handed to the model -- real
public cross-lingual KGs and operational ICT metadata traces -- before any
claim about LLM-supplied cross-lingual knowledge is warranted.

\section*{CRediT authorship contribution statement}
\textbf{Yuxiang Sun:} Conceptualization, Methodology, Software, Investigation,
Writing -- original draft. \textbf{Nan Yin:} Validation, Data curation,
Writing -- review \& editing.

\section*{Declaration of competing interest}
The authors declare that they have no known competing financial interests or
personal relationships that could have appeared to influence the work reported
in this paper.

\section*{Funding}
This work was supported by the 2026 Research Project of Guangzhou University
of Software (grant number KY202649).

\section*{Declaration of generative AI and AI-assisted technologies in the
manuscript preparation process}
During the preparation of this work the authors used an AI assistant in order
to assist with code refactoring, statistical verification and manuscript
editing. After using this tool, the authors reviewed and edited the content as
needed and take full responsibility for the content of the published article.

\section*{Code and data availability}
To support independent verification, the complete audit artefacts are
publicly released under the MIT licence at the repository
\url{https://github.com/CAIST-2026/ict-express-rdf-audit} (release commit
37cab3e): (i) the deterministic benchmark generator with its fixed seeds and
the 500 ground-truth queries; (ii) the alignment-supervision probe; (iii) all
baseline and ablation implementations; (iv) the statistical-analysis
scripts; (v) the per-query result records from which every table, figure,
and number in this letter is derived; and (vi) the Phase~1 matched-supervision,
real-RDF transfer, and exact-CI artefacts, together with a provenance index
mapping each reported value to its source file and script; and (vii) the
systematic-revision package with the unified latency workload, expanded
Wikidata benchmark, consolidated statistics, and reference-verification
records. No real operator,
subscriber, or device data is contained in the released artefacts; the
benchmark is synthetic throughout.

\begin{thebibliography}{00}
\bibitem{pan2024unifying}
S. Pan, L. Luo, Y. Wang, C. Chen, X. Wang, X. Wu, Unifying large language
models and knowledge graphs: a roadmap, IEEE Trans. Knowl. Data Eng. 36 (7)
(2024) 3580--3599. https://doi.org/10.1109/TKDE.2024.3352100.
\bibitem{barnett2024seven}
S. Barnett, S. Kurniawan, S. Thudumu, Z. Brannelly, M. Abdelrazek, Seven
failure points when engineering a retrieval augmented generation system, in:
Proc.\ IEEE/ACM 3rd Int.\ Conf.\ AI Engineering (CAIN), ACM, 2024, pp.\ 194--199.
https://doi.org/10.1145/3644815.3644945.
\bibitem{edge2024graphrag}
D. Edge, H. Trinh, N. Cheng, et al., From local to global: a graph RAG approach
to query-focused summarization, arXiv:2404.16130, 2024.
\bibitem{sun2018jape}
Z. Sun, W. Hu, Q. Zhang, Y. Qu, Bootstrapping entity alignment with knowledge
graph embedding, in: Proc.\ IJCAI, 2018, pp.\ 4396--4402.
https://doi.org/10.24963/ijcai.2018/611.
\bibitem{yang2024segalign}
L. Yang, J. Cheng, F. Zhang, Advancing cross-lingual entity alignment with
large language models, in: Findings of EMNLP 2024, ACL, 2024, pp.\ 8122--8138.
https://doi.org/10.18653/v1/2024.findings-emnlp.475.
\bibitem{zhang2023miracl}
X. Zhang, N. Thakur, O. Ogundepo, et al., MIRACL: a multilingual retrieval
dataset covering 18 diverse languages, Trans. Assoc. Comput. Linguist. 11
(2023) 1114--1131. https://doi.org/10.1162/tacl\_a\_00595.
\bibitem{thorne2018fever}
J. Thorne, A. Vlachos, C. Christodoulopoulos, A. Mittal, FEVER: a large-scale
dataset for fact extraction and VERification, in: Proc.\ NAACL-HLT, 2018,
pp.\ 809--819. https://doi.org/10.18653/v1/N18-1074.
\bibitem{lewis2020rag}
P. Lewis, E. Perez, A. Piktus, et al., Retrieval-augmented generation for
knowledge-intensive NLP tasks, in: Advances in Neural Information Processing
Systems 33 (NeurIPS), 2020, pp.\ 9459--9474.
\bibitem{rdflib}
G. van Harmelen, S. Schalkwijk, et al., rdflib: a Python library for working
with RDF, version 7.6.0, 2025. https://github.com/RDFLib/rdflib.
\bibitem{e5}
X. Wang, et al., Multilingual E5 text embeddings: a technical report,
arXiv:2402.05672, 2024. https://arxiv.org/abs/2402.05672.
\bibitem{qwen25}
A. Yang, et al., Qwen2.5 technical report, arXiv:2412.15115, 2024.
https://arxiv.org/abs/2412.15115.
\bibitem{mcnemar1947}
Q. McNemar, Note on the sampling error of the difference between correlated
proportions or percentages, Psychometrika 12 (2) (1947) 153--157.
\bibitem{holm1979}
S. Holm, A simple sequentially rejective multiple test procedure, Scand. J.
Stat. 6 (2) (1979) 65--70.
\end{thebibliography}

\end{document}
